{  \newpage {\bfseries \fontsize{14}{12} \selectfont \centerline{Abstract} 
\vspace*{2\baselineskip}} \setlength{\parindent}{11mm} 
{ \setlength{\parindent}{0mm} } }

Image Geolocation and Hazard-Aware Response is a multimodal image understanding system that combines place-aware visual reasoning with risk-aware decision support. The system accepts an image from the user, analyzes visual and contextual cues using a vision-capable multimodal large language model, and returns a structured response containing summary, possible location, risk level, confidence, entities, rationale, and recommended actions. For normal landmark or place images, the system provides informational actions such as learn-more and map lookup. For hazardous or suspicious scenes, the system shows report-oriented safety actions only after explicit user confirmation.

The project is implemented as a Python Flask backend, SQLite persistence layer, and React-based web interface. Gemini 2.5 Flash is used as the primary multimodal model, while deterministic post-processing logic validates model output, normalizes confidence, applies hazard keyword rules, and maps risk states to user-facing actions. The system stores sessions, messages, model outputs, actions, and incident records for auditability. The evaluation uses project run logs containing 22 analyzed runs across 14 sessions. Low-risk and high-risk outputs each occurred eight times, uncertain outputs occurred six times, and the average confidence was 0.623. Ten user-confirmed incident records and twelve learn-more actions were logged. The results show that structured multimodal reasoning can support both geospatial understanding and controlled hazard response, while preserving human verification for safety-sensitive actions.
